\section{Discussion and Summary}
\label{sec:summary}

\subsection{Summary of findings}
This assignment studied transfer learning for seven-class dermatoscopic skin-lesion classification on DermaMNIST
(HAM10000, 224$\times$224). Seven controlled experiments led to five conclusions:
\begin{enumerate}
  \item \textbf{Pretraining matters most.} Fine-tuning an ImageNet ResNet-50 lifts test macro-F1 from
        \ScratchFOne\% (CNN trained from scratch) to \RFineFOne\%, and accuracy from \ScratchAcc\% to
        \RFineAcc\%.
  \item \textbf{Fine-tune, do not freeze.} Frozen ImageNet features reach only \RFrozenFOne--\EFrozenFOne\%
        macro-F1. The dermatoscopic domain is far enough from ImageNet that the backbone itself must adapt;
        discriminative learning rates with warm-up make this stable.
  \item \textbf{Gradual unfreezing did not help} in a 12-epoch budget (\RGradFOne\% vs \RFineFOne\%), although it
        is cheaper and ranks well (AUC \RGradAUC).
  \item \textbf{Class weighting trades accuracy and calibration for minority-class recall.} Balanced accuracy goes
        up 4.5 points, \texttt{df}/\texttt{vasc} F1 up 5--6 points, and accuracy down 0.8 points. ECE worsens from
        \RNoWECE{} to \RFineECE.
  \item \textbf{8-view TTA is a reliable, free-to-train gain.} It improves accuracy and AUC for every model.
\end{enumerate}

The final model, \textbf{ResNet-50 with full fine-tuning and 8-view TTA}, selected on validation only, reaches
\textbf{\RFineAccTTA\% accuracy}, \textbf{\RFineBalAccTTA\% balanced accuracy}, \textbf{\RFineFOneTTA\% macro-F1}
(95\% CI \RFineFOneTTACI) and \textbf{macro-AUC \RFineAUCTTA} on the held-out test set. That is about 10 accuracy
points above the best result in the official MedMNIST v2 benchmark table. It detects \BestRecMel\% of melanomas,
and its errors are concentrated in the clinically hard \texttt{mel}/\texttt{nv} and
\texttt{akiec}/\texttt{bkl} pairs. Grad-CAM shows that its evidence usually lies on the lesion itself.

\subsection{Limitations}
\label{sec:limitations}
\begin{itemize}
  \item \textbf{Same-lesion leakage.} The official split is image-level. Our audit (Section~\ref{sec:leakage})
        found re-photographed training lesions in the test set, so the absolute scores are probably slightly
        optimistic. A lesion-level split needs the HAM10000 metadata, which DermaMNIST does not ship.
  \item \textbf{Single seed and single split.} Each configuration was trained once (seed 42). The bootstrap CIs
        capture test-set sampling noise but not training variance. Several of the closer comparisons
        (ResNet-50 vs EfficientNet-B0, weighted vs unweighted) overlap within their CIs and should be repeated
        over 3--5 seeds before drawing firm conclusions.
  \item \textbf{Limited tuning budget.} Learning rates, epochs and the unfreezing schedule were set from common
        practice, not searched, because of the 6\,GB GPU. Gradual unfreezing in particular might do better with a
        longer schedule.
  \item \textbf{Under-confident probabilities.} The final model's raw probabilities are under-confident (ECE
        \RFineECETTA) and should be temperature-scaled before being read as risk scores.
  \item \textbf{Dataset scope.} HAM10000 comes from two centres and mostly fair skin types. Performance on darker
        skin, on other dermatoscopes, or on clinical (non-dermatoscopic) photographs is untested.
\end{itemize}

\subsection{Ethical considerations}
This is a course project on a public research dataset (CC BY-NC 4.0). It is \emph{not} a medical device and must
not be used for diagnosis. In any clinical use, the most costly error is a missed melanoma, so recall on the
malignant classes and calibrated uncertainty matter more than overall accuracy. The model would also need
validation across skin types to avoid unequal performance between patient groups.

\subsection{Future work}
\begin{itemize}
  \item Re-split by lesion using HAM10000 metadata and re-evaluate, to remove the leakage.
  \item Repeat every run over 5 seeds and report mean $\pm$ standard deviation; tune per-class decision thresholds
        and the temperature on the validation split.
  \item Add hair-removal and vignette-aware preprocessing, which Grad-CAM suggests is needed.
  \item Compare stronger backbones (ConvNeXt, Vision Transformers, dermatology-specific foundation models) and
        larger input sizes using the same pipeline~\cite{medmnistplus}.
  \item Test external generalisation on ISIC 2019/2020 images.
\end{itemize}

\subsection{Conclusion}
For a small, imbalanced medical-imaging dataset, ImageNet transfer learning with \emph{full} fine-tuning and
discriminative learning rates turned a weak from-scratch baseline (\ScratchFOne\% macro-F1) into a strong
classifier (\RFineFOneTTA\% macro-F1, \RFineAccTTA\% accuracy with TTA). This was done on a single 6\,GB consumer
GPU in about an hour of total training. The comparison also shows that the details matter: whether to freeze
or fine-tune, how to handle class imbalance, and how to evaluate (macro-F1, confidence intervals, calibration,
leakage checks). These choices change both the numbers and the conclusions one should draw from them.
